\Needspace{8\baselineskip}\originalchapter{2003 作业六: 箍缩与发电}
题源: 21 Oct 2003, due 30 Oct.
\begin{exercise}\label{ex:03-6-1}
原题1. \ensuremath{\theta} pinch平衡 $B_z=B_0+(B_a-B_0)r/a$,0\ensuremath{\le}r\ensuremath{\le}a,p(a)=0. 求 (a)p(r); (b)J(r); (c)\ensuremath{\beta}=$2\mu_0\langle p\rangle/B_a^2$最大值, 平均为 $2a^{-2}\int_0^ap r\dd r$.
\end{exercise}
\begin{solution}
总压力常数, $p=(B_a^2-B_z^2)/(2\mu_0)$, $J_\theta=-(B_a-B_0)/(\mu_0a)$. 令x=B0/Ba, 按r权重积分得 $\beta=1/2-x/3-x^2/6$. 若允许B反转且p\ensuremath{\ge}0, \ensuremath{-}1\ensuremath{\le}x\ensuremath{\le}1, 最大为2/3在x=\ensuremath{-}1. 若附加不反转0\ensuremath{\le}x\ensuremath{\le}1, 最大1/2在x=0. 原题未写不反转, 必须报告这两个不同约束. B在轴处线性r的导数也提示理想剖面轴光滑性有限, 不把该端点模型当严格三维光滑磁场.
\end{solution}
\begin{exercise}\label{ex:03-6-2}
原题2. 纯z pinch无Bz, $J_z=j_0(1-r/a)$, p(a)=0. 求 (a)B(r); (b)p(r); (c)证明p(0)=$\mu_0j_0^2a^2/18$.
\end{exercise}
\begin{solution}
安培圆积分得 $B_\theta=\mu_0j_0[r/2-r^2/(3a)]$. 由p\ensuremath{{}^{\prime}}=\ensuremath{-}JB, 令s=r/a, 直接积分
\begin{equation}
 p=\mu_0j_0^2a^2\left(\frac1{18}-\frac{s^2}4+\frac{5s^3}{18}-\frac{s^4}{12}\right).
\end{equation}
在$s=1$归零, $s=0$给所求中心值; 导数回代验证符号为向外降压.
\end{solution}
\begin{exercise}\label{ex:03-6-3}
原题3. MHD发电机为宽a(x)、高b(y)矩形通道, 均匀密度速度沿z、均匀By, 两x电极电压差\ensuremath{\phi}、均匀电流大小j. (a)\ensuremath{\eta}\ensuremath{\simeq}0求速度; (b)p(0)=p0求p(z); (c)每长度发电功率; (d)压强力每长度做功率; (e)\ensuremath{\eta}固定非零、V与B固定, 调j求最大发电.
\end{exercise}
\begin{solution}
指定取电E沿+x, $E=|\phi|/a$, V>0沿z, 电流 $\mathbf J=-j\mathbf e_x$, j\ensuremath{\ge}0. Ohm给E=VB\ensuremath{-}\ensuremath{\eta}j. (a)V=E/B. (b)磁力\ensuremath{-}jB沿z, 稳态0=\ensuremath{-}p\ensuremath{{}^{\prime}}\ensuremath{-}jB, p=p0\ensuremath{-}jBz. (c)输出 $P_\ell=-ab\mathbf J\cdot\mathbf E=abjE$. (d)压力做功 $abV(-p')=abVjB$, 理想时等于输出; 非理想差为ab\ensuremath{\eta}j\ensuremath{{}^{2}}耗热. (e) $P_\ell=abj(VB-\eta j)$, 最优j=VB/(2\ensuremath{\eta}), $P_{\ell,\max}=abV^2B^2/(4\eta)$. 固定V意味着外供流动功率, 不能让无限大j产生无限功率而保持压强与速度.
\end{solution}
\begin{exercise}\label{ex:03-6-4}
原题4. 曲线的密切平面由切向与曲率半径方向张成. 证明无力MHD平衡的\ensuremath{\nabla}B在磁力线密切平面内.
\end{exercise}
\begin{solution}
J\ensuremath{\times}B=0给 $(\mathbf B\cdot\nabla)\mathbf B=\nabla(B^2/2)$. 写B=Bb, 左为 $B(\dd B/\dd s)\mathbf b+B^2\boldsymbol\kappa$, 因此 $\nabla B=(\dd B/\dd s)\mathbf b+B\boldsymbol\kappa$. 它是切向b和曲率\ensuremath{\kappa}的线性组合. 在直线段$\kappa=0$时密切平面不唯一, 结论为梯度沿切向, 仍无binormal分量.
\end{solution}
\begin{exercise}\label{ex:03-6-5}
原题5. z pinch p(a)=0, 将MHD平衡再积分证明平均p仅由总电流I决定并求函数. 氢n=$10^{20}$m$^{-3}$均匀, Te=Ti=$T_0(1-r^2/a^2)$, T0=10keV,a=.01m, 求I.
\end{exercise}
\begin{solution}
平衡 $p'=-(B/\mu_0r)(rB)'$. 乘r\ensuremath{{}^{2}}积分并分部, p边界零给 $2\int_0^apr\dd r=a^2B_a^2/(2\mu_0)$, 所以 $\langle p\rangle=B_a^2/(2\mu_0)=\mu_0I^2/(8\pi^2a^2)$. 两物种p=2nT0(1\ensuremath{-}r\ensuremath{{}^{2}}/a\ensuremath{{}^{2}}), 平均nT0. 得 $I=\sqrt{8\pi^2a^2nT_0/\mu_0}=3.173\times10^4$A. 这是Bennett类型积分关系, 不保证所选剖面稳定.
\end{solution}
